\bookchapter{第五十二章：四大組件與 Activity 生命之期}

\chapterintro{若謂 Compose 為 App 之「顏」，則此章所言者，App 之「骸」也。安卓系統有四大基本組件，凡 App 皆由之拼成；而其中至要者曰 Activity，有一「生命之期圖」，凡安卓開發者皆須銘之於心。此章既學，則「何切後臺而返，吾 App 之狀態遂失」此類千古之謎，汝能解之矣。}

\booksection{一、四大組件：安卓之「四梁八柱」}

安卓官家定四種基本組件，各司一類之事：

\begin{center}
\begin{tabularx}{\linewidth}{|X|X|X|}
\hline
\textbf{組件} & \textbf{職守} & \textbf{設譬} \tabularnewline
\hline
Activity & 可見之界面 & 店鋪之門面——用戶可見、可交互 \tabularnewline
Service & 後臺任事之「隱者」 & 倉庫之工——奏樂、下載文件，毋須界面 \tabularnewline
BroadcastReceiver & 受系統之廣播 & 傳達之室——「電池無電矣」「網斷矣」之喇叭 \tabularnewline
ContentProvider & 跨 App 共享數據 & 公家資料之櫃——譬如以聯絡人共享於他 App \tabularnewline
\hline
\end{tabularx}
\end{center}

\begin{center}\small\color{gray}安卓四大組件速覽\end{center}

新手之階，汝百分之九十五之光陰皆與 Activity 相交接；Service 與 BroadcastReceiver 者，作音樂之播、消息之推時用之；ContentProvider 用之至寡（系統已封常用之數據矣）。又當慎之：四大組件皆須於 AndroidManifest.xml 中登記戶口，不登則系統莫能尋之。

\booksection{二、Activity 生命之期：安卓開發之首張「藏寶圖」}

機上之 App 隨時可被斷：電話至矣、用戶切而回微信、系統記憶體不足而殺汝……Activity 之一生，即於此圖上循環而遊：

onCreate（出生：初始化）

onStart（可見了）

onResume（前臺交互中）

onPause（部分被遮住）

onStop（完全看不見）

onDestroy（銷燬）

onRestart / onStart / onResume：返前臺時循原路而返

\begin{center}\small\color{gray}圖 4：Activity 生命之期（行之方：上排生 下排離，返則原路而折）\end{center}

\begin{center}
\begin{tabularx}{\linewidth}{|X|X|X|}
\hline
\textbf{回調} & \textbf{何時而發} & \textbf{典型之用} \tabularnewline
\hline
onCreate & Activity 創建之時（一遭） & 載入界面、初化數據 \tabularnewline
onStart & 界面轉而可見 & 始行輕量之邏輯 \tabularnewline
onResume & 得焦點、可交互 & 始動畫、復其播 \tabularnewline
onPause & 見遮其半/切走之頃 & 停動畫、存其稿 \tabularnewline
onStop & 全不可見 & 釋重資、停下載 \tabularnewline
onDestroy & 毀之前 & 終局之清理 \tabularnewline
\hline
\end{tabularx}
\end{center}

\begin{center}\small\color{gray}六生命之期回調之典型用\end{center}

\textbf{避坑之警：轉屏將「殺」Activity}\par

新手至經典之覆車：App 轉屏一遭，輸入框中之字盡失。其故：轉屏則常毀而重建 Activity（歷 onDestroy onCreate 一遭）。解之之法：用 rememberSaveable 存其狀態（Compose 中代 remember 者），重建之時自尋回其值。第六、七章將屢用之。

\booksection{三、Intent：App 內外通行之「車票」}

Intent（意）者，組件之間通信之「車票」也。汝欲自一頁跳於他頁，或啟瀏覽器、起相機，皆須「買票登車」：

\begin{lstlisting}[language=Kotlin]
// 方式一：顯式 Intent——明指所往之頁
val intent = Intent(this, DetailActivity::class.java)
intent.putExtra("name", "張三")     // 車上附"行李"（數據）
startActivity(intent)

// 方式二：隱式 Intent——但言所為，系統自擇其人
val web = Intent(Intent.ACTION_VIEW, Uri.parse("https://developer.android.com"))
startActivity(web)   // 系統啟瀏覽器/或使汝擇
\end{lstlisting}

二法各有其用：顯式者用於己 App 內頁之跳轉；隱式者用於調系統或他 App 之能（啟網頁、撥電話、分享）。Compose 項目之中，於 Composable 內以 LocalContext.current 得 Context，即可 startActivity。

\booksection{四、多頁與導航}

真實之 App 皆有多頁。Compose 官家所薦之導航方略，曰 Navigation Compose（依賴之庫 androidx.navigation）：

\begin{lstlisting}[language=Kotlin]
val navController = rememberNavController()

NavHost(navController, startDestination = "home") {
    composable("home") { HomeScreen(navController) }
    composable("detail/{name}") { entry ->
        val name = entry.arguments?.getString("name")
        DetailScreen(name ?: "未知")
    }
}

// 於頁面中跳轉：
navController.navigate("detail/張三")     // 帶參數
navController.popBackStack()              // 返回上一頁
\end{lstlisting}

\begin{tcbinfo}{小識：頁多則奈何？}
項目既長，宜集諸頁之路由於一文件（譬如 AppNavGraph.kt），統而理之，毋使「跳而何往」之邏輯散於諸頁。此自「能行」趨「可維護」之第一步也。

\end{tcbinfo}

\booksection{五、順而識之：Fragment 為何物（及本書何以不教）}

檢舊之教程，汝必遇 Fragment——其為 Activity 中「可復用界面之片段」，於舊式多頁架構中為絕對之主角。然 Compose 之世，頁面逕以 Composable 函數 + Navigation 路由足矣，官家亦已弱 Fragment 之位。本書故不列之於主線。

貽汝之議：新項目勿主動用 Fragment；然讀舊碼、面試見問之時，知其為「Activity 中所嵌之界面片段」足矣。為學之資，恆先擇 2023 年以後者——安卓技術之棧，其「老黃曆」之新，速於所想。

\begin{tcbpractice}{來動手：為計數器裝「記憶」}
改第五十一章之計數器：(1) 易 remember 為 rememberSaveable，轉屏而驗其數尚在；(2) 加「關於」一鈕，以 Navigation 跳於第二頁而顯版本之號；(3) 於兩頁之 onCreate/onPause 中印日誌（Log.d），以 Logcat 觀生命之期調用之序。三步既畢，生命之期遂非紙上談兵矣。

\end{tcbpractice}

\begin{tcbsummary}{本章小結}
\begin{itemize}
  \item 四大組件：Activity（門面）、Service（後臺）、BroadcastReceiver（喇叭）、ContentProvider（資料櫃）
  \item 生命之期六回調：onCreate onStart onResume onPause onStop onDestroy
  \item 轉屏則毀而重建 Activity：用 rememberSaveable 保其狀態
  \item Intent 者「車票」也：顯式跳己頁，隱式調系統之能
  \item 多頁用 Navigation Compose：NavHost + composable 路由
\end{itemize}

\end{tcbsummary}
